\chapter{State-Space Models and the Kalman Filter}\label{ch:qm:state-space-models-and-the-kalman-filter}

A desk hedges Brent crude with WTI and sets the hedge ratio by a 60-day rolling regression of Brent's daily price changes on WTI's. The ratio
has not been stable: fitted year by year it went from 0.41 in 2014 to 0.94 in 2021 and 1.01 in 2026. When the relationship shifts, the rolling
regression takes a month to cover half the move. A \omterm{def:qm:state-space-models-and-the-kalman-filter:kf}{Kalman filter} that treats the ratio as a hidden state drifting from day to day promises to
follow shifts in days. Fitted by maximum likelihood to sixteen years of prices, the filter's answer is sobering: the ratio drifts so slowly
relative to the day's noise that the optimal filter is itself a month-long average, and it hedges 0.7\% better than the window. A faster filter
is available, at a price the likelihood refuses to pay. This chapter is about \omterm{def:qm:state-space-models-and-the-kalman-filter:ssm}{state-space models} and the filter that estimates their hidden
states: the linear Gaussian model and the Kalman recursions, the likelihood they deliver, the smoother and the EM algorithm, \omterm{def:qm:state-space-models-and-the-kalman-filter:pf}{particle filters} for
the models that are not linear or not Gaussian, and the time-varying hedge ratio that tests them.

\section{Linear Gaussian state-space models}

\begin{definition}[State-space model, local level model]\label{def:qm:state-space-models-and-the-kalman-filter:ssm}
A linear Gaussian \emph{state-space model}\index{state-space model} has an observation equation $y_t = Z_t\alpha_t + \varepsilon_t$, $\varepsilon_t \sim \mathcal N(0,
H)$, and a state equation $\alpha_{t+1} = T\alpha_t + \eta_t$, $\eta_t \sim \mathcal N(0, Q)$, with independent noises and $\alpha_1 \sim \mathcal N(a_1, P_1)$. The
\emph{local level model}\index{local level model} is the scalar case $Z = T = 1$: an observed series equals a random-walk level plus noise.
\end{definition}

The form covers a great deal: ARMA models (chapter~17) with the lagged values as the state, a regression whose coefficients drift ($Z_t = x_t^\top$,
$T = I$), a term structure whose factors move over time, a fair value observed through noisy prints. For the hedge, $y_t$ is Brent's daily price change,
$Z_t$ is WTI's, and the state is the hedge ratio $\beta_t$, a random walk: $\Delta b_t = \beta_t\Delta w_t + \varepsilon_t$, $\beta_{t+1} = \beta_t + \eta_t$.

\begin{definition}[Signal-to-noise ratio]\label{def:qm:state-space-models-and-the-kalman-filter:snr}
In a \omterm{def:qm:state-space-models-and-the-kalman-filter:ssm}{local level model}, or a random-walk coefficient model, the \emph{signal-to-noise ratio}\index{signal-to-noise ratio} is $q = Q/H$, the variance of the
state's daily move over the variance of the observation noise.
\end{definition}

\section{The Kalman filter}

\begin{definition}[Kalman filter, innovation, Kalman gain]\label{def:qm:state-space-models-and-the-kalman-filter:kf}
The \emph{Kalman filter}\index{Kalman filter} (Kalman, 1960) computes recursively the conditional law $\alpha_t \mid y_1, \dots, y_{t-1} \sim \mathcal N(a_t, P_t)$.
Given $(a_t, P_t)$, the \emph{innovation}\index{innovation} is $v_t = y_t - Z_ta_t$, with variance $F_t = Z_tP_tZ_t^\top + H$; the \emph{Kalman gain}\index{Kalman gain}
$K_t = P_tZ_t^\top F_t^{-1}$ updates the state, $a_{t|t} = a_t + K_tv_t$ and $P_{t|t} = P_t - K_tZ_tP_t$, and the prediction step gives $a_{t+1} = Ta_{t|t}$ and
$P_{t+1} = TP_{t|t}T^\top + Q$.
\end{definition}

\begin{proposition}[The Kalman filter is Gaussian conditioning]\label{prop:qm:state-space-models-and-the-kalman-filter:kf}
If $\alpha_t \mid \mathcal F_{t-1} \sim \mathcal N(a_t, P_t)$, then $\alpha_t \mid \mathcal F_t \sim \mathcal N(a_{t|t}, P_{t|t})$ with the update above, and
$\alpha_{t+1} \mid \mathcal F_t \sim \mathcal N(a_{t+1}, P_{t+1})$.
\end{proposition}

\begin{proof}
Given $\mathcal F_{t-1}$, $(\alpha_t, y_t)$ is jointly normal with means $(a_t, Z_ta_t)$, variances $P_t$ and $F_t$, and covariance $P_tZ_t^\top$. The
conditional law of a normal vector given one of its components has mean $a_t + P_tZ_t^\top F_t^{-1}(y_t - Z_ta_t)$ and variance $P_t - P_tZ_t^\top
F_t^{-1}Z_tP_t$. The prediction step applies the linear state equation with independent noise.
\end{proof}

It is the Bayesian update of chapter~14 run once a day, with yesterday's posterior as today's prior. A missing observation is handled by skipping the
update: the prediction step alone widens $P$. That is how the chapter treats 20 April 2020, when the EIA spot price of WTI settled at $-36.98$ dollars:
the WTI changes of that day and the next ($-55.29$ and $+45.89$ dollars, against Brent's $-2.39$ and $-8.24$) are not a hedge relationship but a storage
crisis (commercial inventories at Cushing, Oklahoma, had risen by 27 million barrels since mid-March to 83\% of working capacity, and the WTI futures
contract traded below zero for the first time since trading began in 1983, according to the EIA), and both days are set to missing.

\begin{proposition}[Steady-state gain of the local level model]\label{prop:qm:state-space-models-and-the-kalman-filter:gain}
For the \omterm{def:qm:state-space-models-and-the-kalman-filter:ssm}{local level model} with \omterm{def:qm:state-space-models-and-the-kalman-filter:snr}{signal-to-noise ratio} $q$, the predicted variance converges to $P = pH$ with $p = \frac12(q + \sqrt{q^2 + 4q})$, and the gain to
$K = p/(1 + p)$. The filtered level is then an exponentially weighted average of past observations with weight $K$ on the newest.
\end{proposition}

\begin{proof}
In steady state $P = P - P^2/(P + H) + Q$, that is $P^2 = Q(P + H)$; dividing by $H^2$, $p^2 - qp - q = 0$. The recursion $a_{t+1} = a_t + K(y_t - a_t) =
(1 - K)a_t + Ky_t$ is an exponential average.
\end{proof}

For small $q$ the gain is about $\sqrt q$, so the filter's memory is about $1/\sqrt q$ days and its \omterm{def:qm:stochastic-differential-equations:ou}{half-life} $\ln2/\sqrt q$
(\cref{fig:qm:state-space-models-and-the-kalman-filter:gain}). The \omterm{def:qm:state-space-models-and-the-kalman-filter:snr}{signal-to-noise ratio}, not the analyst, decides how fast the filter reacts. For
a regression coefficient the same holds with $q$ multiplied by the regressor's mean square.

\begin{omfigure}
\begin{tikzpicture}
\begin{semilogxaxis}[width=10cm,height=5cm,xlabel={signal-to-noise ratio $q = Q/H$},ylabel={steady-state gain $K$},
  tick label style={font=\small},label style={font=\small},xmin=1e-4,xmax=100,ymin=0,ymax=1,grid=major,grid style={gray!25},table/col sep=comma]
\addplot[omDef,very thick] table[x=q,y=gain]{figdata/methods/19-state-space-models-and-the-kalman-filter/gain.csv};
\addplot[black,dashed,thin,domain=1e-4:0.3,samples=40] {sqrt(x)};
\end{semilogxaxis}
\end{tikzpicture}

\omcaption{Steady-state \omterm{def:qm:state-space-models-and-the-kalman-filter:kf}{Kalman gain} of the \omterm{def:qm:state-space-models-and-the-kalman-filter:ssm}{local level model} against the \omterm{def:qm:state-space-models-and-the-kalman-filter:snr}{signal-to-noise ratio}, $K = p/(1 + p)$ with $p = (q + \sqrt{q^2 + 4q})/2$, and its
small-$q$ approximation $\sqrt q$ (dashed). At $q = 0.01$ the filter puts 9.5\% weight on each new observation; at $q = 1$, 62\%. Data: the chapter's
formula.}\label{fig:qm:state-space-models-and-the-kalman-filter:gain}
\end{omfigure}

\begin{definition}[Prediction-error decomposition]\label{def:qm:state-space-models-and-the-kalman-filter:ped}
The \emph{prediction-error decomposition}\index{prediction-error decomposition} writes the likelihood of the observations as the product of the one-step
predictive densities, $\ln L = -\frac12\sum_t\bigl(\ln(2\pi F_t) + v_t^2/F_t\bigr)$, with the \omterm{def:qm:state-space-models-and-the-kalman-filter:kf}{innovations} and their variances from the filter.
\end{definition}

The filter therefore delivers the likelihood of $(H, Q)$ at the cost of one pass, and maximum likelihood fits the noise variances. On the hedge, one more
refinement matters: daily oil price changes are far from homoskedastic (the standard deviation of WTI's was 0.78 dollars in 2017 and 3.62 in 2026), and a
constant $H$ lets the volatile years dictate the fit. The chapter takes $\Var(\varepsilon_t) = H s_t^2$ with $s_t$ an exponentially weighted volatility of
WTI's changes known the day before, which is the same as filtering the scaled equation $\Delta b_t/s_t = \beta_t\Delta w_t/s_t + \varepsilon_t/s_t$.

\section{A time-varying hedge ratio}

Fitted to the 4\,097 daily changes from January 2010 to September 2026, the scaled model gives $H = 0.508$, $Q = 0.000238$ and a \omterm{def:qm:state-space-models-and-the-kalman-filter:snr}{signal-to-noise ratio} of
0.00047; with the scaled regressor's mean square of 1.12, the filter's half-response to a shift is $\ln2/\sqrt{0.00047 \times 1.12} = 30$ days, exactly the
30 days a 60-day window needs to cover half a shift. The filtered ratio's uncertainty is itself large: its standard deviation is about 0.10, so the daily
ratio is known only to $\pm 0.2$ (\cref{fig:qm:state-space-models-and-the-kalman-filter:ratio}).

\begin{omfigure}
\begin{tikzpicture}
\begin{axis}[width=12cm,height=5.4cm,xlabel={year},ylabel={hedge ratio (Brent on WTI)},
  tick label style={font=\small},label style={font=\small},xmin=2019,xmax=2026.8,ymin=0,ymax=1.6,grid=major,grid style={gray!25},
  xticklabel style={/pgf/number format/set thousands separator={}},
  legend columns=4,legend style={font=\footnotesize,at={(0.5,-0.3)},anchor=north,draw=none,fill=none,
  /tikz/every even column/.append style={column sep=6pt}},table/col sep=comma]
\addplot[name path=hi,draw=none,forget plot] table[x=year,y=hi]{figdata/methods/19-state-space-models-and-the-kalman-filter/ratio.csv};
\addplot[name path=lo,draw=none,forget plot] table[x=year,y=lo]{figdata/methods/19-state-space-models-and-the-kalman-filter/ratio.csv};
\addplot[omDef!15] fill between[of=hi and lo];
\addplot[omDef,thick] table[x=year,y=kalman]{figdata/methods/19-state-space-models-and-the-kalman-filter/ratio.csv};
\addplot[black,very thick] table[x=year,y=smooth]{figdata/methods/19-state-space-models-and-the-kalman-filter/ratio.csv};
\addplot[omThm,thin] table[x=year,y=roll]{figdata/methods/19-state-space-models-and-the-kalman-filter/ratio.csv};
\legend{$\pm 2$ sd,Kalman filter,smoother,60-day regression}
\end{axis}
\end{tikzpicture}

\omcaption{The Brent-on-WTI hedge ratio of daily price changes, 2019--2026: the Kalman-filtered ratio with two standard deviations of its uncertainty, the
smoothed ratio (which uses the whole sample), and the 60-day rolling regression. Data: FRED series DCOILBRENTEU and DCOILWTICO (US Energy Information
Administration spot prices).}\label{fig:qm:state-space-models-and-the-kalman-filter:ratio}
\end{omfigure}

Each estimate hedges day $t$ with the ratio known at the close of day $t - 1$. Over the 4\,035 days after a 60-day start, the mean squared hedge error, in
dollars squared per barrel per day, is 3.324 unhedged, 1.329 with the ratio of the first 60 days held fixed, 1.276 with the 60-day regression and 1.267 with
the \omterm{def:qm:state-space-models-and-the-kalman-filter:kf}{Kalman filter}: a reduction of 0.7\% against the window. Without the volatility scaling, the maximum-likelihood filter is 1.4\% \emph{worse} than the
window, because the volatile years push its $Q$ up tenfold. Making the filter faster is possible: multiplying $Q$ by 10 cuts the half-response to 9.6 days and
raises the hedge error by 0.6\% relative to the window; multiplying it by 100 gives 3.0 days and 6.2\% more error
(\cref{fig:qm:state-space-models-and-the-kalman-filter:tradeoff}). In a simulation where the ratio jumps from 0.7 to 1.0 with the fitted noise levels, the
likelihood's filter and the window both need about fifty days to reach 0.85; the filter with 100 times the $Q$ needs eight.

\begin{omfigure}
\begin{tikzpicture}
\begin{semilogxaxis}[width=10cm,height=5cm,xlabel={days to cover half a shift in the ratio},ylabel={relative hedge error},
  tick label style={font=\small},label style={font=\small},xmin=2,xmax=150,ymin=0.97,ymax=1.08,grid=major,grid style={gray!25},
  xtick={3,10,30,100},xticklabels={3,10,30,100},
  yticklabel style={/pgf/number format/fixed,/pgf/number format/precision=2},
  legend style={font=\footnotesize,at={(0.97,0.97)},anchor=north east,fill=white,draw=none},table/col sep=comma]
\addplot[omDef,very thick,mark=*,mark size=1.8pt] table[x=halfdays,y=relvar]{figdata/methods/19-state-space-models-and-the-kalman-filter/tradeoff.csv};
\addplot[only marks,mark=square*,mark size=2.5pt,omThm] coordinates {(30,1.0)};
\addplot[only marks,mark=o,mark size=4pt,black,thick] coordinates {(30.26,0.9928)};
\legend{Kalman filter ($Q \times 0.1$ to $\times 100$),60-day regression,maximum likelihood}
\end{semilogxaxis}
\end{tikzpicture}

\omcaption{Speed against accuracy for the Brent--WTI hedge, 2010--2026: each point is a \omterm{def:qm:state-space-models-and-the-kalman-filter:kf}{Kalman filter} with the fitted $Q$ scaled by 0.1, 0.3, 1, 3, 10, 30 and
100, placed by its half-response time and its mean squared hedge error relative to the 60-day regression. The likelihood's choice (circled) sits beside the
window; a filter that follows shifts in days pays for it in noise. Data: FRED series DCOILBRENTEU and DCOILWTICO (EIA).}\label{fig:qm:state-space-models-and-the-kalman-filter:tradeoff}
\end{omfigure}

\section{The smoother and expectation--maximisation}

\begin{definition}[Kalman smoother]\label{def:qm:state-space-models-and-the-kalman-filter:smoother}
The \emph{Kalman smoother}\index{Kalman smoother} computes $\E[\alpha_t \mid y_1, \dots, y_n]$ and its variance from the filter's output by a backward pass
(Rauch, Tung and Striebel, 1965): $\hat\alpha_t = a_{t|t} + J_t(\hat\alpha_{t+1} - a_{t+1})$ with $J_t = P_{t|t}T^\top P_{t+1}^{-1}$.
\end{definition}

The smoother answers a different question from the filter: not what the desk could have known, but what the ratio most likely was. It is the tool for
research on past data (\cref{fig:qm:state-space-models-and-the-kalman-filter:ratio} shows it smoother and earlier than the filter) and a source of look-ahead
bias when it leaks into a backtest.

\begin{definition}[Expectation--maximisation algorithm]\label{def:qm:state-space-models-and-the-kalman-filter:em}
The \emph{expectation--maximisation algorithm}\index{expectation--maximisation algorithm} (Dempster, Laird and Rubin, 1977) maximises a likelihood with
unobserved variables by alternating an E-step, the expected complete-data log-likelihood given the data and the current parameters, and an M-step that
maximises it. For \omterm{def:qm:state-space-models-and-the-kalman-filter:ssm}{state-space models} the E-step is the smoother and the M-step for $H$ and $Q$ has a closed form (Shumway and Stoffer, 1982).
\end{definition}

Each EM step cannot decrease the likelihood: the complete-data log-likelihood's expectation is a lower bound on the log-likelihood, up to a constant, that
touches it at the current parameters, and the M-step raises the bound. EM is stable and needs no derivatives, but it can be slow. Started from $H = 1$ and $Q =
0.01$ on the hedge model, it raises the log-likelihood from $-4\,947$ to $-4\,524$ in ten steps and to $-4\,502$ in forty, still 32 points short of the maximum
$-4\,470$ that direct maximisation finds (\cref{fig:qm:state-space-models-and-the-kalman-filter:em}): the likelihood is flat along the direction that trades $Q$
against $H$.

\begin{omfigure}
\begin{tikzpicture}
\begin{axis}[width=10cm,height=4.8cm,xlabel={EM iteration},ylabel={log-likelihood},
  tick label style={font=\small},label style={font=\small},xmin=0,xmax=41,ymin=-4600,ymax=-4450,grid=major,grid style={gray!25},
  yticklabel style={/pgf/number format/set thousands separator={\,}},
  legend style={font=\footnotesize,at={(0.97,0.05)},anchor=south east,fill=white,draw=none},table/col sep=comma]
\addplot[omDef,very thick,mark=*,mark size=1pt] table[x=iter,y=loglik]{figdata/methods/19-state-space-models-and-the-kalman-filter/em.csv};
\addplot[black,dashed,thick] table[x=iter,y=mle]{figdata/methods/19-state-space-models-and-the-kalman-filter/em.csv};
\legend{{EM from $H = 1$, $Q = 0.01$},maximum likelihood}
\end{axis}
\end{tikzpicture}

\omcaption{EM on the scaled Brent--WTI hedge model: the log-likelihood rises at every step (from $-4\,947$ at the start, off the chart) but approaches the maximum
found by direct optimisation (dashed, $-4\,470$) slowly. Data: FRED series DCOILBRENTEU and DCOILWTICO (EIA).}\label{fig:qm:state-space-models-and-the-kalman-filter:em}
\end{omfigure}

\section{Particle filters}

\begin{definition}[Particle filter, sequential importance resampling]\label{def:qm:state-space-models-and-the-kalman-filter:pf}
A \emph{particle filter}\index{particle filter} represents the filtering law of a state by a weighted sample of particles. \emph{Sequential importance
resampling}\index{sequential importance resampling}, the bootstrap filter of Gordon, Salmond and Smith (1993), propagates each particle through the state
equation, weights it by the density of the new observation given the particle, and resamples in proportion to the weights; the average weight at each step
estimates the one-step predictive density, and their product the likelihood.
\end{definition}

The \omterm{def:qm:state-space-models-and-the-kalman-filter:pf}{particle filter} needs neither linearity nor Gaussian noise, only the ability to simulate the state and evaluate the observation density; its price is
Monte Carlo error and degeneracy, measured by the \omterm{def:qm:bayesian-methods:ess}{effective sample size} of chapter~14 computed from the weights, $1/\sum_iw_i^2$. On a linear Gaussian
local level of 500 observations, where the \omterm{def:qm:state-space-models-and-the-kalman-filter:kf}{Kalman filter} is exact, 2\,000 particles give a log-likelihood of $-785.06$ against the exact $-784.22$ and filtered
means within 0.07 of the exact ones (whose posterior standard deviation is 0.45), with an average \omterm{def:qm:bayesian-methods:ess}{effective sample size} of 1\,678.

The useful case is nonlinear. A stochastic-volatility model for daily EUR/USD returns lets the log variance follow an AR(1), $x_t = \mu + \phi(x_{t-1} - \mu)
+ \sigma_\eta u_t$, and observes $r_t = e^{x_t/2}z_t$: no \omterm{def:qm:state-space-models-and-the-kalman-filter:kf}{Kalman filter} applies. With 5\,000 particles on the last 1\,000 returns, a grid search of the particle
likelihood (same random numbers at every point) picks $\phi = 0.9$ and $\sigma_\eta = 0.4$ (log-likelihood $-568.4$, against $-572.3$ at $\phi = 0.95$,
$\sigma_\eta = 0.2$: a flat surface). The filtered volatility follows the GARCH volatility of chapter~18 but reacts faster
(\cref{fig:qm:state-space-models-and-the-kalman-filter:sv}); the \omterm{def:qm:bayesian-methods:ess}{effective sample size} averages 4\,449 but falls to 23 on 3 April 2025, when the euro rose 2.7\%
after the US tariff announcement of 2 April (a risk-off episode in which, unusually, the dollar depreciated, as the ECB's Financial Stability Review of November 2025
notes) and almost every particle was too calm to explain it.

\begin{omfigure}
\begin{tikzpicture}
\begin{axis}[width=12cm,height=5cm,xlabel={year},ylabel={daily volatility (\%)},
  tick label style={font=\small},label style={font=\small},xmin=2022.8,xmax=2026.8,ymin=0,ymax=1.5,grid=major,grid style={gray!25},
  xtick={2023,2024,2025,2026},xticklabel style={/pgf/number format/set thousands separator={}},
  legend columns=2,legend style={font=\footnotesize,at={(0.5,-0.3)},anchor=north,draw=none,fill=none,
  /tikz/every even column/.append style={column sep=8pt}},table/col sep=comma]
\addplot[omDef,thick] table[x=year,y=sv]{figdata/methods/19-state-space-models-and-the-kalman-filter/sv.csv};
\addplot[omThm,very thick,dashed] table[x=year,y=garch]{figdata/methods/19-state-space-models-and-the-kalman-filter/sv.csv};
\legend{{stochastic volatility, particle filter},GARCH-$t$ (chapter~18)}
\end{axis}
\end{tikzpicture}

\omcaption{EUR/USD daily volatility over the last 1\,000 ECB fixings to 23 September 2026: the particle-filtered stochastic-volatility estimate ($\phi = 0.9$,
$\sigma_\eta = 0.4$, 5\,000 particles) and the GARCH-$t$ conditional volatility. Data: ECB euro reference rates (source: ECB
statistics).}\label{fig:qm:state-space-models-and-the-kalman-filter:sv}
\end{omfigure}

\section{Tutorial: the hedge ratio that moved}\label{tut:qm:state-space-models-and-the-kalman-filter:hedge}

\begin{tutorial}
\textbf{Goal.} Track the Brent--WTI hedge ratio with a \omterm{def:qm:state-space-models-and-the-kalman-filter:kf}{Kalman filter} fitted by maximum likelihood, and measure what it buys against a rolling regression.
\textbf{End state:} \cref{fig:qm:state-space-models-and-the-kalman-filter:ratio,fig:qm:state-space-models-and-the-kalman-filter:tradeoff}, the \omterm{def:qm:state-space-models-and-the-kalman-filter:snr}{signal-to-noise
ratio} 0.00047 and the 0.7\% reduction in hedge error.
\begin{tutsteps}
\item \textbf{The filter}, with missing observations and the prediction-error likelihood.
\omcode{code/firm/kalman/firm_kalman.py}{32}{58}{The Kalman filter with the prediction-error decomposition.}\label{lst:qm:state-space-models-and-the-kalman-filter:kf}
\item \textbf{The streaming tracker} the desk would run each evening.
\omcode{code/firm/kalman/firm_kalman.py}{174}{188}{A streaming time-varying hedge ratio.}\label{lst:qm:state-space-models-and-the-kalman-filter:tracker}
\item \textbf{Run} \texttt{compare()}, \texttt{tradeoff()}, \texttt{em\_path()} and \texttt{sv\_grid()} in \texttt{qm\_kalman.py}, then \texttt{fig\_kalman.py}.
\end{tutsteps}
\textbf{What to change next.} Let the ratio mean-revert ($T = \phi < 1$ with a mean) instead of wandering, and compare likelihoods; add WTI's lagged change as a
second regressor to absorb the time difference between the two spot assessments.
\end{tutorial}

\section{Build: the Kalman toolkit}\label{bld:qm:state-space-models-and-the-kalman-filter:kalman}

\begin{build}
\bfield{Purpose} Hidden quantities the miniature firm tracks every day (hedge ratios, fair values seen through noisy prints, drifting signal loadings) are
filtered here, with their uncertainty.
\bfield{Interface} \texttt{kalman\_filter(y, Z, T, H, Q, a0, P0)}; \texttt{rts\_smoother(kf, T)}; \texttt{em(\dots)}; \texttt{fit\_mle(\dots)};
\texttt{local\_level\_gain(q)}; \texttt{particle\_filter(y, n, init, step, loglik\_obs, rng)}; \texttt{HedgeTracker(q, h, beta0, p0).update(x, y)}.
\bfield{Rules} Filters for trading use only data up to the previous close; smoothers are for research and never feed a backtest; missing data are NaN, not zeros;
every \omterm{def:qm:state-space-models-and-the-kalman-filter:pf}{particle filter} reports its \omterm{def:qm:bayesian-methods:ess}{effective sample size}.
\bfield{Acceptance tests} \texttt{code/firm/kalman/tests/}: the likelihood and the smoothed states equal direct Gaussian conditioning on a short series; a missing
observation skips the update; the steady-state gain formula; maximum likelihood and EM (monotone) recover a simulated local level; the \omterm{def:qm:state-space-models-and-the-kalman-filter:pf}{particle filter}
matches the \omterm{def:qm:state-space-models-and-the-kalman-filter:kf}{Kalman filter} on a linear model; the streaming tracker reproduces the filter.
\bfield{Stretch} Square-root filtering for numerical stability; the extended and unscented filters; particle marginal Metropolis--Hastings for the
stochastic-volatility parameters.
\end{build}

\begin{omsources}
\item EIA, \emph{Today in Energy}, on 2020 crude oil prices (negative WTI on 20 April 2020); ECB, \emph{Financial Stability Review}, November 2025, special feature on the April 2025 tariff announcement.
\item R.~E.~Kalman, ``A new approach to linear filtering and prediction problems'', \emph{Journal of Basic Engineering} 82, 1960.
\item H.~E.~Rauch, F.~Tung and C.~T.~Striebel, ``Maximum likelihood estimates of linear dynamic systems'', \emph{AIAA Journal} 3, 1965.
\item A.~P.~Dempster, N.~M.~Laird and D.~B.~Rubin, ``Maximum likelihood from incomplete data via the EM algorithm'', \emph{Journal of the Royal
Statistical Society B} 39, 1977.
\item R.~H.~Shumway and D.~S.~Stoffer, ``An approach to time series smoothing and forecasting using the EM algorithm'', \emph{Journal of Time Series
Analysis} 3, 1982.
\item N.~J.~Gordon, D.~J.~Salmond and A.~F.~M.~Smith, ``Novel approach to nonlinear/non-Gaussian Bayesian state estimation'', \emph{IEE Proceedings F} 140,
1993.
\item US Energy Information Administration, WTI and Brent spot prices, via FRED (series DCOILWTICO and DCOILBRENTEU), accessed 24 September 2026.
\end{omsources}

\section{Exercises}

\begin{exercise}[$\star$]\label{exo:qm:state-space-models-and-the-kalman-filter:1}
A \omterm{def:qm:state-space-models-and-the-kalman-filter:ssm}{local level model} has $H = 1$ and $Q = 0.04$. What are the steady-state gain and the \omterm{def:qm:stochastic-differential-equations:ou}{half-life} of the filter's memory?
\end{exercise}

\begin{exercise}[$\star$]\label{exo:qm:state-space-models-and-the-kalman-filter:2}
Yesterday's filtered hedge ratio is $0.80$ with variance $0.01$; $Q = 0.0002$, $H = 0.5$. Today WTI moves $+2$ and Brent $+1.2$ (scaled units). Update the ratio.
\end{exercise}

\begin{exercise}[$\star$]\label{exo:qm:state-space-models-and-the-kalman-filter:3}
Why does a missing observation increase the state's variance, and by how much in the \omterm{def:qm:state-space-models-and-the-kalman-filter:ssm}{local level model}?
\end{exercise}

\begin{exercise}[$\star\star$]\label{exo:qm:state-space-models-and-the-kalman-filter:4}
Show that for small $q$ the local level gain is approximately $\sqrt q$, and compare with an exponential moving average: which $\lambda$ does the filter
imply for $q = 0.0004$?
\end{exercise}

\begin{exercise}[$\star\star$]\label{exo:qm:state-space-models-and-the-kalman-filter:5}
Write the AR(1) plus noise, $y_t = x_t + \varepsilon_t$ with $x_t = \phi x_{t-1} + \eta_t$, in state-space form and give its steady-state predicted variance
equation.
\end{exercise}

\begin{exercise}[$\star\star$]\label{exo:qm:state-space-models-and-the-kalman-filter:6}
A \omterm{def:qm:state-space-models-and-the-kalman-filter:pf}{particle filter}'s normalised weights are 0.5, 0.3, 0.1, 0.05 and 0.05. What is the \omterm{def:qm:bayesian-methods:ess}{effective sample size}?
\end{exercise}

\begin{exercise}[$\star\star\star$]\label{exo:qm:state-space-models-and-the-kalman-filter:7}
\emph{Coding.} On the Brent--WTI data, scale the fitted $Q$ by 0.1, 1, 10 and 100 and tabulate the hedge-error variance relative to the 60-day regression and
the half-response time.
\end{exercise}

\begin{exercise}[$\star\star\star$]\label{exo:qm:state-space-models-and-the-kalman-filter:8}
\emph{Find the flaw.} ``We backtested a hedge using the Kalman-smoothed ratio and cut the hedge error by 15\% against the rolling regression.''
\end{exercise}

\section{Problem: The Hedge Ratio That Moved}

\begin{problem}[{Weekend problem --- how fast should a hedge ratio move?}]\label{pb:qm:state-space-models-and-the-kalman-filter:1}
The data are the EIA's daily WTI and Brent spot prices from FRED, 4 January 2010 to 22 September 2026 (4\,097 daily changes). Brent's daily price change is
hedged with WTI's; the hedge ratio for day $t$ must be known at the close of day $t - 1$.

\medskip\noindent\textbf{Part I --- The data.}
\begin{enumerate}
\item What happened on 20 April 2020, and how does the chapter treat it?
\item How did the yearly regression ratio move between 2014, 2021 and 2026?
\item Why scale the observation noise by WTI's volatility?
\item What are the unhedged and the fixed-ratio hedge-error variances?
\item What does the 60-day regression achieve?
\end{enumerate}

\medskip\noindent\textbf{Part II --- The filter.}
\begin{enumerate}[resume]
\item What $H$, $Q$ and \omterm{def:qm:state-space-models-and-the-kalman-filter:snr}{signal-to-noise ratio} does maximum likelihood give?
\item What is the filter's half-response to a shift, and how does it compare with the window's?
\item What hedge-error variance does the filter achieve, and what reduction against the window?
\item What happens without the volatility scaling?
\item How precisely is the daily ratio known?
\end{enumerate}

\medskip\noindent\textbf{Part III --- Speed and noise.}
\begin{enumerate}[resume]
\item What do filters with $Q$ multiplied by 10 and by 100 achieve?
\item In the simulated jump from 0.7 to 1.0, how many days do the filters and the window need to reach 0.85?
\item What does EM find from a poor start, and why is it slow?
\item What are the filtered and rolling ratios at the end of the sample?
\item When would a fast filter be worth its noise?
\end{enumerate}

\medskip\noindent\textbf{Part IV --- Judgement.}
\begin{enumerate}[resume]
\item Should the desk replace its rolling regression?
\item What else could improve the hedge more than the choice of filter?
\item Why must the smoother stay out of the backtest?
\item State the \emph{named result}: the reduction in hedge-error variance and the estimated \omterm{def:qm:state-space-models-and-the-kalman-filter:snr}{signal-to-noise ratio}.
\item In one sentence: what decides how fast a \omterm{def:qm:state-space-models-and-the-kalman-filter:kf}{Kalman filter} reacts?
\end{enumerate}
\end{problem}

\section{Interview questions}

\begin{interviewq}[$\star$ \iqroles{researcher, trader}]\label{iq:qm:state-space-models-and-the-kalman-filter:1}
Explain the \omterm{def:qm:state-space-models-and-the-kalman-filter:kf}{Kalman filter} to a trader in three sentences.
\end{interviewq}

\begin{interviewq}[$\star\star$ \iqroles{researcher}]\label{iq:qm:state-space-models-and-the-kalman-filter:2}
Derive the Kalman update for a scalar state from Gaussian conditioning.
\end{interviewq}

\begin{interviewq}[$\star\star$ \iqroles{researcher, trader}]\label{iq:qm:state-space-models-and-the-kalman-filter:3}
How do you choose the process noise of a \omterm{def:qm:state-space-models-and-the-kalman-filter:kf}{Kalman filter} tracking a hedge ratio?
\end{interviewq}

\begin{interviewq}[$\star\star$ \iqroles{researcher, mle}]\label{iq:qm:state-space-models-and-the-kalman-filter:4}
What is the difference between filtering and smoothing, and why does it matter in backtests?
\end{interviewq}

\begin{interviewq}[$\star\star$ \iqroles{mle, researcher}]\label{iq:qm:state-space-models-and-the-kalman-filter:5}
When would you use a \omterm{def:qm:state-space-models-and-the-kalman-filter:pf}{particle filter} instead of a \omterm{def:qm:state-space-models-and-the-kalman-filter:kf}{Kalman filter}, and what can go wrong?
\end{interviewq}

\begin{interviewq}[$\star\star\star$ \iqroles{researcher}]\label{iq:qm:state-space-models-and-the-kalman-filter:6}
Why does each EM iteration not decrease the likelihood?
\end{interviewq}
